% ============================================================
% audit_table.tex
% Updated 2026-04-29 (post v0.96-RC + Appendix B + P9 concordance)
% ============================================================

\begin{table}[htbp]
\centering
{\small\caption{System Audit: $\Tdiamond$ predictions mapped to
experiments and supporting derivations.  Status semantics in the
post-abstract release notes (\texttt{PASS} / \texttt{PART} /
\texttt{STUB} / \texttt{FAIL}).  \emph{Scoring note:}
\texttt{exp\_12}'s four-significant-figure \texttt{PASS} is scored
on a $k$-scan resonance peak (short-run quantitative match);
long-horizon two-body stability is captured separately by the
\texttt{exp\_12b} and \texttt{exp\_12c} entries below.}}
\label{tab:audit}
\footnotesize
\setlength{\tabcolsep}{4pt}
\renewcommand{\arraystretch}{0.95}
\begin{tabular}{@{}p{2.6cm}p{2.8cm}p{3.2cm}p{5.4cm}p{1.0cm}@{}}
\hline
\textbf{Physical Phenomenon} & \textbf{Lattice Logic} & \textbf{Continuum Limit} & \textbf{Experiment} & \textbf{Status} \\
\hline
% ── Kinematics ───────────────────────────────────────────────────────────────
Speed of light limit
    & Topological adjacency $v_\text{max}=1$
    & $c = \ell_P / t_P$
    & \texttt{exp\_00}
    & \texttt{PASS} \\
Inertial persistence
    & Phase gradient $\nabla\phi$
    & Newton's first law
    & \texttt{exp\_01}
    & \texttt{PASS} \\
$1/r^2$ discrete correction
    & Path count deviation from CLT Gaussian
    & Falsifiable short-range correction
    & \texttt{exp\_06}
    & \texttt{PASS} \\
Dirac equation
    & Bipartite tick rule, $\mathcal{A}=1$
    & $(i\gamma^\mu\partial_\mu - m)\psi = 0$
    & Derived analytically (§\ref{sec:kinematics})
    & \texttt{PASS} \\
Spin-$\tfrac{1}{2}$
    & Two-component spinor from bipartite structure
    & Clifford algebra on $O_h$-averaged frame condition
    & Derived analytically (§\ref{sec:kinematics})
    & \texttt{PASS} \\
Equivalence principle
    & Gravitational and inertial mass both $= \sin(\omega)/2$
    & $m_g = m_i$ structural identity
    & \texttt{exp\_01}, \texttt{exp\_02}
    & \texttt{PASS} \\
Lattice birefringence (kinematic)
    & $M_\text{eff}$ eigenvalues $\{4/3, 4/3, 0\}$;
      flat dispersion along $(1,1,-1)$
    & Direction-resolved photon dispersion (P7 refinement)
    & Derived analytically (§\ref{subsec:dirac_limit}, eq:M\_eff)
    & \texttt{STUB} \\
Zitterbewegung source (dual hypothesis)
    & Bipartite $\omega$-rate alternation: self-driven vs.\ recoil-driven
    & Direction-dependent rest mass under recoil hypothesis (P9)
    & Derived analytically (§\ref{subsec:bipartite_velocity})
    & \texttt{STUB} \\
\hline
% ── Quantum phenomena ─────────────────────────────────────────────────────────
Quantum interference
    & Discrete path summation
    & Schr\"{o}dinger equation
    & \texttt{exp\_03}
    & \texttt{PASS} \\
Decoherence
    & Phase scrambling via observer clock
    & Density matrix decay
    & \texttt{exp\_04}
    & \texttt{PASS} \\
Measurement / irreversibility
    & Observer as clock addition
    & Decoherence theory
    & \texttt{exp\_05}
    & \texttt{PASS} \\
Born rule
    & $\mathcal{A}=1$ conservation $\Rightarrow$ $|\psi|^2$
    & Measurement postulate derived
    & \texttt{exp\_03}--\texttt{exp\_05}
    & \texttt{PASS} \\
Mass quantization
    & Arnold tongue frequency locking
    & Discrete spectrum without boundary condition
    & \texttt{exp\_09}
    & \texttt{PASS} \\
Photon dispersion
    & Group velocity vs.\ wavenumber $k$
    & Brillouin zone boundary
    & \texttt{exp\_09}
    & \texttt{PASS} \\
\hline
% ── Gravity ───────────────────────────────────────────────────────────────────
Gravitational attraction
    & Clock density gradient, refractive bias
    & Newtonian $1/r^2$
    & \texttt{exp\_02}
    & \texttt{PASS} \\
Gravitational time dilation
    & Scheduler load near mass concentration
    & Schwarzschild metric
    & \texttt{exp\_02}, \texttt{exp\_07}
    & \texttt{PASS} \\
Clock density conservation
    & Continuity equation from tick()
    & Einstein field equations (§\ref{sec:clock_fluid})
    & \texttt{exp\_07}
    & \texttt{PASS} \\
Event horizon + Hawking $T$
    & Scheduler saturation $\rho \to \ell_P^{-3}$;
      boundary surface gravity $\kappa = c^4/(4GM)$
    & No singularity; $T_H = \hbar c^3/(8\pi G M k_B)$
    & Derived analytically (§\ref{subsec:saturation})
    & \texttt{PART} \\
Bekenstein-Hawking entropy
    & Boundary session count, $A_\text{session} = 4\ell_P^2$
    & $S = k_B A/(4\ell_P^2) = k_B A c^3/(4G\hbar)$
    & Derived analytically (§\ref{subsec:black_holes})
    & \texttt{PART} \\
\hline
% ── Electromagnetism ──────────────────────────────────────────────────────────
EM vs.\ gravity (curl vs.\ div)
    & curl($\phi$) vs.\ div($\phi$) phase geometry
    & Maxwell + Einstein
    & \texttt{exp\_08}
    & \texttt{PASS} \\
Photon emission
    & $\mathcal{A}=1$ session creation at orbit lock-in
    & Spontaneous emission rate
    & \texttt{exp\_08}
    & \texttt{PASS} \\
Photon emission as $\mathcal{A}=1$ necessity
    & Joint $\mathcal{A}=1$ + momentum-transfer recoil at orbit;
      tongue-jump mandates third session
    & Spontaneous emission without separate field quantization
    & \texttt{exp\_19c} v10 sweeps 2026-04-26 showed 0 recoils
      across rates $0.01$--$0.5$; \texttt{exp\_20} 2026-05-02
      three-arm comparison (drain, unitary beam splitter, phase
      rotation) did not yet exhibit orbital lock-in -- baseline
      two-body instability dominates (\texttt{exp\_12b}); next
      attempt awaits resolution of the long-horizon stability
      question
    & \texttt{PART} \\
Emission-operator joint $\mathcal{A}=1$ conservation
    & Pointwise unitary 2-mode rotation between electron and photon
      sessions with joint amplitude enforced
    & QED postulates a separately quantised electromagnetic field
    & \texttt{exp\_20} 2026-05-02 sweep at $\text{GRID}=65^3$,
      rate $=0.05$, 6000 ticks per arm.  Joint $\mathcal{A}=1$
      preserved to machine precision ($8.9\!\times\!10^{-16}$) by
      arm B (beam splitter); arm A (the as-written drain in
      \texttt{exp\_19c}) violates by $1.0$; arm C (phase rotation)
      trivial.  Phase-dependent non-monotonic transfer in arm B is
      an unanticipated refinement of the operation-algebra
      prediction (see \nolinkurl{notes/lattice_operation_algebra.md}
      row~6 caveat)
    & \texttt{PART} \\
Induced gauge action (Sakharov form)
    & $-\mathrm{Tr}\ln D_\text{lat}[U]$;
      bipartite plaquette sum dominates at leading order
    & Maxwell density on $O_h$-averaged observables;
      structural form derived (numerical $1/g^2$ open)
    & Derived analytically (Appendix~\ref{sec:induced_gauge_action})
    & \texttt{PART} \\
Gauge-sector birefringence
    & Plaquette-sum tensor $\mathbf{Q}$ eigenvalues $\{4, 4, 16\}$;
      oblate ellipsoid in $\mathbf{B}$-space
    & Direction-dependent photon kinetic term;
      same optical axis $(1,1,-1)$ as kinematic sector (P9)
    & Derived analytically (Appendix~\ref{sec:induced_gauge_action})
    & \texttt{STUB} \\
\hline
% ── Hydrogen and atomic physics ───────────────────────────────────────────────
Hydrogen $E_n \sim -1/n^2$
    & Coulomb well orbital resonance
    & Bohr spectrum
    & \texttt{exp\_10}
    & \texttt{PASS} \\
Spontaneous quantization ($n=1$)
    & Arnold tongue attractor in Coulomb well
    & $k_\text{Bohr} = 1/R_1$
    & \texttt{exp\_11}
    & \texttt{PASS} \\
Two-body hydrogen (4 sig.\ figs.)
    & Live proton session, joint resonance
    & $k_\text{min} = 0.0970$ vs.\ $k_\text{Bohr} = 0.0971$
    & \texttt{exp\_12}
    & \texttt{PASS} \\
Two-body long-horizon stability
    & Bare \texttt{exp\_12} chassis run for 6000 ticks (no emission)
    & Standard QM treats Bohr orbits as stationary
    & \texttt{exp\_12b} 2026-05-02 single run at $\text{GRID}=65^3$;
      $r_\text{peak}$ escapes from $\sim R_1$ to grid edge ($\sim 83$)
      by tick $\sim 2000$; $\mathcal{A}=1$ preserved at machine
      precision ($8.9\!\times\!10^{-16}$).  Refines but does not
      invalidate \texttt{exp\_12}'s 4-sig-fig PASS, which is scored
      on a shorter $k$-scan resonance peak
    & \texttt{PART} \\
Long-horizon escape mechanism
    & Bare two-body orbit metastable; escape onset varies with
      grid size
    & Standard QM has no lattice substrate to exhibit finite-volume
      effects
    & \texttt{exp\_12c} 2026-05-05 grid sweep: orbit escapes at
      tick $\sim 140$ on $\text{GRID} \geq 81^3$ and tick $\sim 2000$
      on $65^3$.  The $65^3$ ``stability'' is a finite-volume
      artefact: boundary reflections confine the wave packet near
      $R_1$.  True orbital lifetime is hundreds of ticks, not
      thousands.  Init-mode and drift-mode controls rule out
      initial-condition asymmetry and systematic operator drift as
      causes
    & \texttt{PART} \\
Lock-in resonance grid-independence
    & \texttt{exp\_12} chassis on $\text{GRID} \in \{65, 81, 97,
      113\}^3$, pre-escape window (tick 30--119)
    & Standard QM has no lattice substrate to vary
    & \texttt{exp\_12d\_tight} 2026-05-05 four-grid sweep:
      $r$ at $k_\text{Bohr} = 0.0971$ in $[11.54, 12.30]$
      across all grids ($\Delta r \approx 0.8$); three of four
      grids' $k_\text{min}$ within $\pm 0.003$ of $k_\text{Bohr}$;
      $r_\text{std} \in [1.1, 2.9]$ across all configurations
      except two isolated outliers at $k=0.101$ on 81 and 113
      traced to early-escape tail of \texttt{exp\_12c}'s
      metastable-orbit mechanism (\texttt{exp\_12d\_outlier\_trace},
      abrupt escape at tick 119); lock-in attractor is
      grid-independent
    & \texttt{PART} \\
Bound state requires active proton
    & Static well collapses; live proton explores well
    & Proton as symmetry-breaking session
    & \texttt{exp\_11} ($n=2$), \texttt{exp\_12}
    & \texttt{PASS} \\
Three-body helium-like system
    & Two-electron + proton sessions
    & Helium ground state
    & \texttt{exp\_13}
    & \texttt{PASS} \\
Two-electron helium
    & Same-sublattice exclusion + joint resonance
    & Helium spectrum
    & \texttt{exp\_14}
    & \texttt{PASS} \\
Proton mass lower bound
    & Minimum $\Omega_P$ for quantization;
      binding $\neq$ quantization
    & All swept $\Omega_P$ produce Regime~2 (bound, unquantized);
      threshold lies above swept range
    & \texttt{exp\_16}
    & \texttt{PART} \\
Standard Model gauge group as automorphism
    & $\mathrm{Aut}(\Tdiamond, \mathcal{A}=1) \stackrel{?}{=}
      SO(3,1) \times SU(3) \times SU(2) \times U(1)$
      on extended per-site $\mathbb{C}^{12}$ amplitude
    & 18-dimensional Lie-algebra equality;
      structure constants must match
      $\mathfrak{so}(3,1) \oplus \mathfrak{su}(3) \oplus
      \mathfrak{su}(2) \oplus \mathfrak{u}(1)$ factor-by-factor
    & Six \nolinkurl{src/utilities/automorphism_*} and
      \nolinkurl{tick_rule_*} sympy scripts (2026-05-05) establish:
      discrete spatial $\mathrm{Aut}(\Gamma, V)$ is order 48
      ($B_3 \cong O_h$ abstractly, 12 elements orthogonal);
      existing RGB symmetry contributes only $\mathbb{Z}_3
      \subset SU(3)$; per-site $SU(2)$ on $\mathbb{C}^2$ overlaps
      with Lorentz rotations (not a separate factor); the
      conjecture's direct-product structure is achievable on the
      extended $\mathbb{C}^{12} = \mathbb{C}^2 \otimes
      \mathbb{C}^2 \otimes \mathbb{C}^3$ amplitude (verified,
      all 18 generators commute pairwise across factors);
      trivial tensor extension preserves $\mathcal{A}=1$ and
      parity (verified in
      \nolinkurl{tick_rule_extended_consistency.py});
      Wilson plaquette $V_1, -V_2, -V_1, V_2$ is gauge-invariant
      (verified in \nolinkurl{tick_rule_gauge_invariance.py}).
      Three structural questions open: exact equality of
      Aut, SM-chirality coupling, explicit $1/g^2$ prefactors
    & \texttt{STUB} \\
\hline
% ── Falsifiability framework ─────────────────────────────────────────────────
Multi-channel concordance (P9)
    & Single optical axis $(1,1,-1)$ inherited by kinematic, gauge,
      gravitational, and recoil-mass sectors
    & Five-channel direction agreement $(a/\lambda)^2$;
      time-dependence variants distinguish self-projecting from
      externally-driven embedding
    & Derived analytically (§\ref{subsec:p9_concordance});
      direction-resolved GRB / CMB / LIGO / atomic-clock data
    & \texttt{STUB} \\
\hline
\end{tabular}
\normalsize
\end{table}
